\chapter{Code reference}
\label{ch:reference}

Paths are relative to \code{code/subroutines}. Each routine has a help text
(\code{help <name>}) that documents its input and output. The call graph is short:
the studies of the paper lift a mesh to degree $p$ (\code{usemBuildModel},
\code{pSpectralElement_NodesTopology}), compute element matrices (\code{Ke_2D} or its
vectorized form \code{cyl_Ke_PM}) from the shape functions, and solve with
\texttt{MATLAB}'s direct solver.

\section{Element technology}

\begin{xltabular}{\linewidth}{@{}P{6.4cm}L@{}}
  \toprule
  routine & purpose \\
  \midrule
  \endhead
  \code{ElementMatrices/Ke_2D(p, C, x, y [, nIPxi, nIPeta])} & unsymmetric element
    stiffness $\int\limits \tilde{\matr B}_\mathrm{L}^{\T}\matr C\matr B_\mathrm{R}$
    with $\tilde{\matr B}_\mathrm{L}=\det(\matr J)\,\matr B_\mathrm{L}$, built
    with $\adj(\matr J)$ (default $2p$ Gauss
    points per direction) \\
  \code{ShapeFunctions/2D/ShapeFunctions2D_SEM(xi, eta, p)} & parametric (test)
    shape functions and $\partial/\partial\xi$, $\partial/\partial\eta$;
    \texttt{[nNodes x nPoints]} arrays \\
  \code{ShapeFunctions/2D/ShapeFunctions2D_Metric(xQ, yQ, xN, yN, p [, method, xiQ])} &
    metric (trial) shape functions and $\partial/\partial x$, $\partial/\partial y$;
    \code{method} 1 = $\mathcal Q_p$ in skew coordinates (default),
    2 = hybrid space (\cref{sec:metric-hybrid}; \code{xiQ}: local coordinates of the
    query points); empty \code{method}: \code{p.metricMethod}; basis of the
    Vandermonde matrix from \code{p.metricBasis} (default \code{'lagrange'},
    \cref{ch:metric-basis}); optional 4th output $\matr V$ \\
  \code{ShapeFunctions/2D/ShapeFunctions2D_MetricHybrid} & hybrid space \\
  \code{ShapeFunctions/2D/CheckUnisolvent(V, px, py)} & error if $\matr V$ is singular
    to working precision \\
  \code{ShapeFunctions/2D/KronRows(A, B)} & column-wise Kronecker product \\
  \code{ShapeFunctions/2D/ReferenceNodes2D(p)} & local coordinates of the element
    nodes in the order of \code{ShapeFunctions2D_SEM} \\
  \code{ShapeFunctions/2D/CorotationalSkewCoordinates(xQ, yQ, xN, yN, p)} & skew
    coordinates $(s,t)$ and $\matr J_0$ \\
  \code{ShapeFunctions/2D/ShapeFunctions2D_4nodes}, \code{_8nodes} & bilinear and
    Serendipity functions (geometry of the input mesh) \\
  \code{ShapeFunctions/1D/LagrangePolynomials(x)} & 1D Lagrange polynomials
    (barycentric form) as function handles; with the helpers
    \code{LagrangeBarycentric}, \code{LagrangeBarycentricDerivative},
    \code{LagrangeNodePolynomial}, \code{LagrangeNodeDerivative} \\
  \code{ShapeFunctions/1D/ShapeFunctions1D_SEM(p, xi, ndist)} & 1D shape
    functions, order \code{[xi=-1, xi=+1, interior]} \\
  \code{ShapeFunctions/1D/PolynomialBasis1D(s, n, basis, ndist)} & 1D bases
    (monomial, Legendre, Chebyshev, Lagrange on the reference nodes) and their
    derivatives, vectorized \\
  \code{GeometryMapping/Global2LocalNewton(xQ, yQ, xN, yN, p)} & inverse mapping
    $\vect x\mapsto(\xi,\eta)$ by Newton's method, vectorized \\
  \code{GeometryMapping/Jacobi_Matrix(p, [xi, eta], x, y)} &
    $\matr J = [x_{,\xi}\ y_{,\xi};\ x_{,\eta}\ y_{,\eta}]$ \\
  \code{GeometryMapping/Local2GlobalCoordinatesQuad(p, P, [x, y])} &
    $\vect x(\xi,\eta)$ \\
  \code{Quadrature/lgwt(n, a, b)} & Gauss-Legendre points and weights
    (G.~von Winckel) \\
  \code{Quadrature/lglnodes(N)} & $N+1$ Gauss-Lobatto-Legendre points and
    weights (G.~von Winckel) \\
  \code{NodalDistribution/nodal_distribution(p, ndist)} & supporting points on
    $[-1,1]$ (\code{equispaced_nodes}, \code{lglnodes}, \code{cgl_nodes}) \\
  \bottomrule
\end{xltabular}

\section{Mesh and topology}

\begin{xltabular}{\linewidth}{@{}P{6.4cm}L@{}}
  \toprule
  routine & purpose \\
  \midrule
  \endhead
  \code{ReadInputFileAbaqus/Read_mesh_ABAQUS(file)} & nodes and 8-node elements of an
    \texttt{ABAQUS} input file \\
  \code{Topology/pSpectralElement_NodesTopology(T, x, y, p, p_org)} & creates the
    edge and interior nodes of degree $p$ from the 4-node ($p_\mathrm{org}=1$) or
    8-node ($p_\mathrm{org}=2$) mesh; shared edges get the same nodes, in flipped
    order for opposite orientation \\
  \code{Topology/CheckCoincidentNodes} & error if an element has coincident nodes \\
  \code{Topology/Topology_Edge_2D(T)} & global edge numbers \texttt{[nElem x 4]}
    (edges 1-2, 2-3, 3-4, 4-1) \\
  \code{Topology/ConnectivityEdgeElement(TE)} & elements and local edge numbers of
    every edge \\
  \code{Topology/Topology_DoF_2D(T)} & DOF connectivity, node $n\mapsto$ DOFs
    $2n-1$, $2n$ \\
  \bottomrule
\end{xltabular}

\section{Material}

\begin{xltabular}{\linewidth}{@{}P{6.4cm}L@{}}
  \toprule
  routine & purpose \\
  \midrule
  \code{Material/HookeMatrix3d(E, nu)} & isotropic $6\times6$ stiffness \\
  \code{Material/PlaneStressConstitutiveMatrices(C)} & static condensation of
    $\sigma_{33}=0$, $3\times3$ matrix \\
  \code{Material/HookeMatrix2d(E, nu, antype)} & closed-form $3\times3$ matrices
    (1 plane stress, 2 plane strain) \\
  \bottomrule
\end{xltabular}

\section{Routines of the studies}

\begin{xltabular}{\linewidth}{@{}P{6.4cm}L@{}}
  \toprule
  routine & purpose \\
  \midrule
  \endhead
  \code{manuscript/studies/ms_Ke}, \code{ms_Me}, \code{ms_assemble_KM} & element stiffness
    and consistent mass of \PP, \PM\ and \PH, global $\matr K$, $\matr M$ with the
    calibrated rules \\
  \code{manuscript/studies/ms_lebesgue} & Lebesgue constants of one element \\
  \code{manuscript/studies/ms_write_table} & writes a table file \\
  \code{studies/thick_cylinder/matlab/cyl_*} & model with exact arcs, element
    stiffness (\code{cyl_Ke_PM}, \code{cyl_Ke_PP}), assembly, loads, Dirichlet
    conditions, errors and diagnostics of the thick cylinder \\
  \code{studies/rajendran_panel/matlab/rp_*} & model, exact solutions, loads and
    supports of the panel \\
  \bottomrule
\end{xltabular}

\section{Third-party code}

\code{lgwt} and \code{lglnodes} are by G.~von Winckel (File Exchange); they keep the
notices in their file headers.
